\problem{}
ShanghaiTech is planning a night security patrol system on campus.
The campus is modeled as an undirected graph $G=(V,E)$, where each vertex represents an area, and an edge between two vertices means that the two areas are directly connected.

If a patrol point is placed in an area, then that patrol point can monitor both the area itself and all areas directly connected to it.

The university would like to place at most $k$ patrol points so that every area on campus is monitored.

Formally, given an undirected graph $G=(V,E)$ and a positive integer $k$, decide whether there exists a subset $S\subseteq V$ such that
\begin{itemize}
	\item $|S|\le k$, and
	\item for every vertex $v\in V$, either $v\in S$, or there exists some $u\in S$ such that $(u,v)\in E$.
\end{itemize}

Prove that this decision problem is NP-complete.
\solution{
\textbf{1. Membership in NP.}
For a proposed set $S$, we check $|S|\le k$ and then check every
vertex to see whether it is in $S$ or has a neighbor in $S$.

\textbf{2. Known NP-complete problem.}
We reduce from \textsc{Vertex Cover}, which is NP-complete.

\textbf{3. Polynomial-time reduction.}
Let $(G,k)$ be a \textsc{Vertex Cover} instance, where $G=(V,E)$. We first
ignore isolated vertices, since they are irrelevant for vertex cover.

Construct a graph $G'=(V',E')$ as follows. Keep all vertices and edges of $G$.
For every edge $e=\{u,v\}\in E$, add a new vertex $x_e$ and two edges
$\{x_e,u\}$ and $\{x_e,v\}$. Thus
\[
V'=V\cup \{x_e:e\in E\}.
\]
The new parameter is still $k$. This construction is polynomial time.

\textbf{4. Correctness of the reduction.}
Now suppose $G$ has a vertex cover $C$ with $|C|\le k$. We show that the same
set $C$ dominates $G'$. For each new vertex $x_e$, at least one endpoint of
$e$ is in $C$, so $x_e$ is dominated. For an original vertex $v$, either
$v\in C$, or $v$ has an incident edge $\{v,u\}$ and the cover property forces
some endpoint of this edge to be in $C$. Since $v\notin C$ in this case, we
must have $u\in C$, so $v$ is dominated by $u$. Thus $C$ is a dominating set of
size at most $k$ in $G'$.

For the other direction, let $D$ be a dominating set of $G'$ with $|D|\le k$.
If $D$ contains a new vertex $x_e$, where $e=\{u,v\}$, replace $x_e$ by one of
its endpoints, say $u$. This does not increase the size and does not destroy
domination: $x_e$ is dominated by $u$, the endpoint $u$ is dominated by itself,
the other endpoint $v$ is dominated by $u$ because $\{u,v\}$ is an original
edge, and $x_e$ has no other neighbors. Repeating this replacement, we obtain a
dominating set $D'\subseteq V$ with $|D'|\le k$.

For any original edge $e=\{u,v\}$, the new vertex $x_e$ must be dominated by
$D'$. Since $D'\subseteq V$ and $x_e$ is adjacent only to $u$ and $v$, at least
one of $u,v$ belongs to $D'$. Therefore every edge of $G$ is covered by $D'$.

So $G$ has a vertex cover of size at most $k$ iff $G'$ has a dominating set of
size at most $k$. Therefore the patrol-point problem is NP-hard. Since it is
also in NP, it is NP-complete.
}

\newpage
